\documentclass[12pt,a4paper]{article}
\usepackage{fontspec}
\usepackage[russian]{babel}
\usepackage{geometry}
\usepackage{booktabs}
\usepackage{longtable}
\usepackage{array}
\usepackage{ragged2e}
\usepackage{hyperref}
\usepackage{enumitem}
\usepackage{amsmath}
\usepackage{amssymb}
\usepackage{listings}

\geometry{left=25mm,right=18mm,top=20mm,bottom=22mm}
\setmainfont{Times New Roman}
\setsansfont{Arial}
\hypersetup{colorlinks=true,linkcolor=black,urlcolor=blue}
\setlist{noitemsep,topsep=3pt}
\renewcommand{\arraystretch}{1.18}
\lstset{basicstyle=\ttfamily\footnotesize,breaklines=true,frame=single,columns=fullflexible}
\newcolumntype{L}[1]{>{\RaggedRight\arraybackslash}p{#1}}
\sloppy
\emergencystretch=4em

\title{Краткое описание пайплайна MNT4 pairing verifier}
\author{Итоговая архитектура проекта}
\date{22 мая 2026}

\begin{document}
\maketitle

\section{Что реализовано}
Первая часть --- полный on-chain reference для MNT4. Он нужен как контрольная реализация и как экспериментальная демонстрация того, что прямое вычисление MNT4-сопряжения в EVM практически непригодно: актуальный тестовый запуск дает $259\,719\,954$ gas для reference-вычисления digest сопряжения.

Вторая часть --- основной вариант работы. Сопряжение считается off-chain в Rust. Solidity-контракт не пересчитывает цикл Миллера и финальную экспоненту. Он получает точки, digest результата, commitments к вспомогательным данным и proof-envelope, затем проверяет, что публичные входы proof согласованы с переданным утверждением.

Общий поток:
\begin{lstlisting}
Rust MNT4 backend
  -> MNT4 artifacts and relation roots
  -> compact BN254 proof envelope
  -> Solidity verifier
\end{lstlisting}

BN254 используется только как дешевый EVM-совместимый слой проверки proof, потому что Ethereum имеет precompile для BN254 pairing check. Это не заменяет MNT4-вычисление: MNT4-арифметика и MNT4-сопряжение строятся в Rust backend и сверяются с arkworks.

\section{Off-chain часть}
\subsection{Rust backend}
Основной backend находится в:
\begin{lstlisting}
offchain_verify/crates/mnt4_trace_backend
\end{lstlisting}

Главные файлы:
\begin{itemize}
    \item \texttt{src/main.rs} --- CLI-вход: читает request и пишет JSON-артефакты;
    \item \texttt{src/lib.rs} --- построение MNT4 artifact, commitments, relation roots и proof input.
\end{itemize}

Backend принимает точки $P_i$, точку $Q$, режим fixed-Q или Parametric-Q, context, epoch и служебные поля. В canonical fixture используются $P$ и $Q$, соответствующие генераторам MNT4-753 из arkworks.

\subsection{Что считается}
Основная функция:
\begin{lstlisting}
build_artifact(request) -> PairingArtifact
\end{lstlisting}

Она выполняет следующие шаги.
\begin{enumerate}
    \item Нормализует точки $P_i$ и $Q$ в формате MNT4-753.
    \item Для $Q$ строит коэффициенты линий цикла Миллера через arkworks \texttt{G2Prepared::from(q)}.
    \item Упаковывает линии в sparse-формат: отдельно линии удвоения и линии сложения.
    \item Считает commitments:
    \[
        C_{dbl}=H(\text{double lines}),\quad
        C_{add}=H(\text{addition lines}),\quad
        C_{lines}=H(C_{dbl},C_{add}).
    \]
    \item Считает Miller loop через \texttt{MNT4\_753::multi\_miller\_loop}.
    \item Выполняет final exponentiation и получает digest результата сопряжения.
    \item Строит relation roots: \texttt{lineCacheRelationRoot}, \texttt{millerRelationRoot}, \texttt{finalExponentiationRelationRoot}.
    \item Записывает summary MNT-native relation fragment в поле \texttt{nativeRelation}.
\end{enumerate}


\subsection{Выходные файлы}
\texttt{write\_artifacts} создает:
\begin{itemize}
    \item \texttt{trace.json} --- общий artifact;
    \item \texttt{witness.json} --- witness-данные;
    \item \texttt{public\_inputs.json} --- публичные входы;
    \item \texttt{proof\_input.json} --- вход для генерации proof-envelope.
\end{itemize}

\section{MNT-native relation fragment}
Отдельный Rust crate:
\begin{lstlisting}
offchain_verify/crates/mnt_cycle_constraints
\end{lstlisting}

Он нужен для продолжения идеи: MNT4/MNT6 образуют цикл, поэтому отношение для будущего folding можно строить в MNT-sized setting, а не эмулировать MNT4 полностью в BN254.

Текущий compiled relation fragment принимает три публичных root:
\begin{itemize}
    \item root кэша линий;
    \item root Miller relation;
    \item root final exponentiation relation.
\end{itemize}

Для canonical fixture получено:
\begin{itemize}
    \item compiled relation constraints: $24\,181$;
    \item public inputs: $3$;
    \item witness variables: $72\,537$;
    \item relation satisfied: true.
\end{itemize}

Это pre-folding слой, который показывает порядок стоимости relation-части до оборачивания в folding.

\section{On-chain контракты}
\footnotesize
\begin{longtable}{L{0.28\textwidth}L{0.62\textwidth}}
\toprule
\textbf{Контракт} & \textbf{Роль} \\
\midrule
\texttt{MNT4PairingFinal} & Главный контракт. В конструкторе создает proof checker и proof-system verifier. Пользователь вызывает именно его. \\
\texttt{MNT4PairingVerifier} & Основная логика проверки: сверяет точки, digest результата, commitments и public inputs proof. \\
\texttt{MNT4PairingProofChecker} & Нейтральный adapter: проверяет statement hash и передает proof в generated verifier. \\
\texttt{MNT4ProofSystemVerifier} & Generated verifier, который использует BN254 precompiles Ethereum. \\
\texttt{MNT4TatePairing}, \texttt{MNT4Extension}, \texttt{BigIntMNT} & Арифметика MNT4 для baseline и тестов. \\
\bottomrule
\end{longtable}
\normalsize

\section{Главный вызов}
Основной метод:
\begin{lstlisting}
verifyPairingClaim(
    G1Point[] pointsP,
    G2Point q,
    bytes32 expectedResultDigest,
    bytes32 artifactRoot,
    bytes32 transcriptHash,
    bytes32 context,
    uint64 epoch,
    ArtifactCommitments commitments,
    bytes proof
) returns (bool)
\end{lstlisting}

Контракт не доверяет переданному кэшу и не принимает digest результата ``на слово''. Он проверяет, что public inputs proof совпадают с утверждением пользователя: hash всех точек $P_i$, hash точки $Q$, digest результата, artifact root, transcript hash и commitments к line-cache, Miller trace и final exponentiation.

Если пользователь меняет хотя бы один объект, проверка падает. После binding-проверок вызывается:
\begin{lstlisting}
PROOF_CHECKER.verify(statementHash, proof)
\end{lstlisting}

\section{Тесты}
\subsection{End-to-end MNT4 verifier}
Файл:
\begin{lstlisting}
offchain_verify/test/final_mnt4_pairing/MNT4PairingFinalE2E.t.sol
\end{lstlisting}

Разворачивает \texttt{MNT4PairingFinal(8)}, берет одну точку $P$ и одну точку $Q$ из canonical fixture, подает digest результата, commitments и proof-envelope. Ожидаемый результат --- \texttt{true}.

В этом же файле есть negative suite: подмена $P$, подмена $Q$, подмена result digest, подмена line commitment, подмена Miller trace commitment и подмена final exponentiation commitment. Во всех случаях контракт должен отклонить вызов.

\subsection{Сравнение MNT4 с arkworks}
Файл:
\begin{lstlisting}
offchain_verify/test/final_mnt4_pairing/MNT4FinalArkworksCrossCheck.t.sol
\end{lstlisting}

Проверяется, что Solidity reference-арифметика согласована с canonical vectors из arkworks:
\begin{itemize}
    \item операции в $F_p$;
    \item операции в $F_{p^2}$;
    \item операции в $F_{p^4}$;
    \item digest MNT4-сопряжения на тех же fixture-точках $P$ и $Q$.
\end{itemize}

\subsection{Fuzz-тест boolean-семантики verifier-ов}
Тест находится в:
\begin{lstlisting}
offchain_verify/test/final_mnt4_pairing/PairingVerifierBooleanSemanticsFuzz.t.sol
\end{lstlisting}

Для одного fuzz-сценария accept/reject оба verifier-пути возвращают одинаковый boolean. Корректность MNT4-арифметики на самих MNT4-точках проверяется отдельно через arkworks-векторы.

Foundry генерирует параметр 	exttt{uint8 scenario}. Если 	exttt{scenario} четный, выбирается валидный сценарий; если нечетный, выбирается невалидный сценарий.
\begin{enumerate}
    \item Валидный сценарий. MNT4-путь получает canonical fixture $P,Q$, корректный digest результата, commitments и proof-envelope. Он возвращает \texttt{true}. BN254-путь вызывает precompile \texttt{0x08} на валидной pairing equation $e(O,G_2)=1$, где $O=(0,0)$ --- point at infinity в EIP-197 encoding, а $G_2$ --- стандартный BN254 G2 generator. Он тоже возвращает \texttt{true}.
    \item Невалидный сценарий. MNT4-путь получает тот же fixture, но первая limb координаты $x$ точки $P$ изменяется через \texttt{x[0] \textasciicircum= 1}; контракт ловит ошибку binding и результат приводится к \texttt{false}. BN254-путь получает одну пару $(G_1,G_2)$, для которой проверка $e(G_1,G_2)=1$ ложна; precompile возвращает \texttt{false}.
\end{enumerate}

Используемые BN254 точки:
\begin{itemize}
    \item $G_1=(1,2)$;
    \item $O=(0,0)$;
    \item $G_2$ в EIP-197 order: \texttt{x\_real=0x198e...12c2}, \texttt{x\_imag=0x1800...f6ed}, \texttt{y\_real=0x0906...75b}, \texttt{y\_imag=0x12c8...7daa}.
\end{itemize}

Текущий результат: fuzz-тест проходит на 256 запусках. Gas по последнему запуску: среднее значение --- $287\,398$ gas, медиана --- $424\,280$ gas.


\subsection{Proof checker tests}
Файл:
\begin{lstlisting}
offchain_verify/test/final_mnt4_pairing/MNT4PairingProofChecker.t.sol
\end{lstlisting}

Проверяется, что валидный proof fixture проходит через checker, statement hash связан с public signal, а подмена public signal отклоняется.

\subsection{Rust-тесты}
Rust backend проверяется отдельно:
\begin{itemize}
    \item \texttt{mnt4\_trace\_backend}: 15 тестов для line-cache relation, Miller relation, final exponentiation relation, tamper-cases и записи \texttt{nativeRelation};
    \item \texttt{mnt\_cycle\_constraints}: 10 тестов для operation model, compiled R1CS fragment и negative tamper test для witness root.
\end{itemize}

\section{Основные результаты}
\footnotesize
\begin{longtable}{L{0.67\textwidth}r}
\toprule
\textbf{Метрика} & \textbf{Значение} \\
\midrule
Полное on-chain reference-вычисление MNT4 pairing digest & 259,719,954 gas \\
Основной on-chain verifier call & 343,435 gas \\
Backend-only proof checker & 325,283 gas \\
Fuzz boolean semantics, mean & 287,398 gas \\
Fuzz boolean semantics, median & 424,280 gas \\
Rust backend total generation & 809 ms \\
Proof generation & 526 ms \\
BN254 verifier-envelope & 1,538 constraints \\
MNT-native accounting model & 24,178 constraints \\
Compiled MNT-native relation fragment & 24,181 constraints \\
Compiled native public inputs & 3 \\
Compiled native witness variables & 72,537 \\
BN254 emulated pairing reference & 1,393,318 constraints \\
Sonobe-like decider reference & 9,000,000 constraints \\
\bottomrule
\end{longtable}
\normalsize

\section{Оценка будущего CycleFold}
Сам CycleFold в этой работе не реализован. Реализован слой, который должен стать его входом:
\[
    \text{MNT4 artifacts} \rightarrow \text{MNT-native relation fragment} \rightarrow \text{future folding}.
\]

Для одной пары текущий compiled relation fragment стоит:
\[
    C_{rel}(1)=24\,181\ \text{constraints}.
\]

Для multi-pairing используется shared Miller accumulator. По текущей модели:
\[
    C_{Miller}(1)=4\,514,\quad
    C_{Miller}(2)=7\,520,\quad
    C_{Miller}(4)=13\,532.
\]

Оценка relation-части без overhead самого folding:
\[
    C_{rel}(n) \approx C_{lines}+C_{Miller}(n)+C_{FE}+3,
\]
где $C_{lines}=19\,554$, $C_{FE}=110$.

\begin{longtable}{L{0.40\textwidth}r}
\toprule
\textbf{Число пар} & \textbf{Relation constraints} \\
\midrule
$n=1$ & 24,181 \\
$n=2$ & 27,187 \\
$n=4$ & 33,199 \\
\bottomrule
\end{longtable}

Полная стоимость будущего folding-step будет иметь вид:
\[
    C_{fold-step}(n)=C_{rel}(n)+C_{fold-overhead}.
\]

Текущая работа не доказывает готовый CycleFold, но показывает главное предварительное условие: MNT-native relation-часть находится на уровне десятков тысяч constraints. Поэтому следующий этап должен folding-ить именно этот MNT-native слой, а не полный non-native MNT4 pairing circuit внутри BN254.

On-chain стоимость будущей версии при сохранении BN254 compression-envelope ожидается того же порядка, что сейчас:
\[
    G_{verify} \approx 3.2\cdot 10^5 - 4.0\cdot 10^5\ \text{gas},
\]
если контракт не пересчитывает MNT4-арифметику on-chain, а проверяет компактный proof.

\end{document}
